\documentclass[aspectratio=169,11pt,t,xcolor=table]{beamer}
\usepackage{slidestyle}
\externaldocument[H-]{handout}
% Lesson 2 lost-mark points (from the material plus teaching observation), numbered once
\setcounter{lostmark}{0}
\lostmark{noy}{Stopping at $x$: a \tturn is a point, so find $y$ from $f(x)$ as well.}
\lostmark{nojust}{Writing ``maximum'' with no reason. Justify with $f''(x)$ or the gradient on each side.}
\lostmark{interval}{Increasing or decreasing: giving the \tturns instead of an interval, or the inequality the wrong way round.}
\lostmark{twovars}{Differentiating a formula that still has two variables. Use the constraint to get one variable first.}
\lostmark{wrongq}{Answering the wrong quantity: the question asks for the maximum area, not the $x$ that gives it.}
\lostmark{context}{No units or context in the final sentence (m, s, \$, people per hour).}

\newlength\figunit
\begin{document}

% ------------------------------------------------------------------
\begin{frame}{Lesson 2 \textperiodcentered\ Turning Points and Optimisation}\relax
\vskip6mm
{\sEmph\bfseries Learning intention\par}\vskip2mm
Use $f'(x)=0$ to find \tturns, decide maximum or minimum, find where a function increases or decreases, and solve maximum and minimum problems.
\vskip6mm
{\sEmph\bfseries Success looks like\par}\vskip2mm
\begin{itemize}
\item I can find a \tturn and justify its nature.
\item I can give an interval where $f$ is increasing or decreasing.
\item I can set up and solve an optimisation problem, with units.
\end{itemize}
\note{Teaching line, not a question slide.\par
Timing: Do Now 10 min, three rounds about 20 min each, practice sheet 30 min, wrap-up 5 min.}
\end{frame}

% ------------------------------------------------------------------
\begin{frame}{Do Now}\relax
\begin{columns}[T]
\begin{column}{0.48\textwidth}
\qn{1} Differentiate $y=(x-3)(x+5)$.
\vskip5mm
\qn{2} Find the gradient of $f(x)=x^{3}-5x$ at the point where $x=2$.
\end{column}
\begin{column}{0.48\textwidth}
\qn{3} Solve $3x^{2}-12=0$.
\vskip5mm
\qn{4} Solve $2x^{2}-x-6=0$.
\end{column}
\end{columns}
\sfoot{Lesson 1 and prerequisites}
\note{Q1 $\dfrac{dy}{dx}=2x+2$\par
Q2 $7$\par
Q3 $x=2$ or $x=-2$\par
Q4 $x=2$ or $x=-1.5$\par
Short solution:\par
Q1 expand: $y=x^{2}+2x-15$\par
Q2 $f'(x)=3x^{2}-5$; $f'(2)=12-5$\par
Q3 $x^{2}=4$\par
Q4 $(2x+3)(x-2)=0$\par
Watch for:\par
Q3 missing the negative root\par
Q4 sign of the $-1.5$ root\par
Diagnostic:\par
Q1 or Q2 wrong: Lesson 1 not secure; re-teach the power rule before round 1\par
Q3 wrong: expect missing turning points today (both roots of $f'(x)=0$)\par
Q4 wrong: let the student use EQUA for $f'(x)=0$ today}
\end{frame}

% ------------------------------------------------------------------
\begin{frame}{The one idea}\relax
\begin{columns}[c]
\begin{column}{0.48\textwidth}
At a \tturn the \ttangent is flat:
\fbx{f'(x)=0}
\vskip2mm
Maximum: gradient $+$ then $-$.\\[1mm]
Minimum: gradient $-$ then $+$.
\end{column}
\begin{column}{0.50\textwidth}\centering
\setlength\figunit{12mm}
% f(x) = x^3 - 6x^2 + 9x + 1: maximum (1,5), minimum (3,1).  Flat tangents drawn at both;
% gradient signs from f'(x) = 3(x-1)(x-3): f'(0)=9>0, f'(2)=-3<0, f'(4)=9>0.  Caller sets \figunit.
\begin{tikzpicture}[x=\figunit, y=0.55\figunit, line cap=round]
  \draw[-{Stealth[length=2mm]}, line width=0.6pt] (-0.6,0) -- (5.6,0) node[right, inner sep=1pt] {$x$};
  \draw[-{Stealth[length=2mm]}, line width=0.6pt] (0,-0.5) -- (0,7.3) node[above, inner sep=1pt] {$y$};
  \begin{scope}
    \clip (-0.6,-0.5) rectangle (4.8,7.2);
    \curve{-0.25}{4.45}{\x*\x*\x-6*\x*\x+9*\x+1}
  \end{scope}
  \draw[line width=0.8pt, dashed] (0.4,5) -- (1.6,5);
  \draw[line width=0.8pt, dashed] (2.4,1) -- (3.6,1);
  \fill (1,5) circle (2.2pt); \fill (3,1) circle (2.2pt);
  \node[above, inner sep=2pt] at (1,5.1) {maximum};
  % 'minimum' right of the flat tangent: at x = 3.7 the curve is at 2.8, above the label
  \node[anchor=west, inner sep=2pt] at (3.65,1) {minimum};
  % gradient signs, placed beside the curve
  \node at (0.6,2.2) {\textbf{$+$}};   % below the rising curve (f(0.6) = 4.5), clear of the y-axis
  \node at (2.35,4.3) {\textbf{$-$}};
  \node at (4.35,3.1) {\textbf{$+$}};
\end{tikzpicture}

\end{column}
\end{columns}
\sfoot{\hp{idea}}
\note{Teaching line, not a question slide.\par
Trace the curve with a finger: up, flat, down, flat, up.\par
Ask: what is the gradient exactly at the top? Zero, because the tangent is horizontal.}
\end{frame}

% ------------------------------------------------------------------
\begin{frame}{Finding turning points}\relax
{\sEmph\bfseries Write it in this order\par}\vskip1mm
\begin{enumerate}
\item Differentiate.
\item Set $f'(x)=0$; solve for $x$.
\item Find each $y$ from $f(x)$.
\item Maximum or minimum? Say why.
\end{enumerate}
\vskip3mm
\qn{1} Find the \tturns of $f(x)=x^{3}-6x^{2}+9x+1$ and state their nature.
\sfoot{\hp{find}}
\note{Q1 maximum $(1,\,5)$; minimum $(3,\,1)$\par
Short solution:\par
Q1 $f'(x)=3x^{2}-12x+9=3(x-1)(x-3)$; $x=1$ or $x=3$\par
Q1 $f(1)=5$, $f(3)=1$\par
Q1 $f''(x)=6x-12$: $f''(1)=-6<0$ max; $f''(3)=6>0$ min\par
Watch for:\par
stopping at $x=1$, $x=3$ (L1)\par
$y$-values from $f'(x)$ give $0$, a sign of the wrong function}
\end{frame}

% ------------------------------------------------------------------
\begin{frame}{Maximum or minimum?}\relax
\renewcommand{\arraystretch}{1.4}
\begin{tabular}{@{}p{50mm}p{55mm}p{45mm}@{}}
\toprule
test & what you see & conclusion\\
\midrule
gradient each side & $+\ \to\ 0\ \to\ -$ & \\
 & $-\ \to\ 0\ \to\ +$ & \\
second derivative & $f''(x)<0$ & \\
 & $f''(x)>0$ & \\
\bottomrule
\end{tabular}
\vskip4mm
$f''(x)$: differentiate twice. Always \kw{write the reason} in a sentence.
\sfoot{\hp{nature}}
\note{Row 1 maximum\par
Row 2 minimum\par
Row 3 maximum, bends down\par
Row 4 minimum, bends up\par
Short solution:\par
sentence frame: $f''(1)=-6<0$, so the turning point at $(1,\,5)$ is a maximum\par
Watch for:\par
L2: a bare “maximum” with no test is not a justification\par
$f''=0$: the test fails; use the gradient either side instead}
\end{frame}

% ------------------------------------------------------------------
\begin{frame}{Practice}\relax
\qn{1} Find the \tturn of $y=x^{2}-8x+3$ and state whether it is a maximum or a minimum.
\vskip5mm
\qn{2} Find the \tturns of $f(x)=2x^{3}-3x^{2}-12x$ and justify their nature.
\vskip5mm
\qn{3} The curve $y=x^{3}+kx$ has a \tturn at $x=2$. Find $k$.
\sfoot{\hp{find}}
\note{Q1 $(4,\,-13)$, minimum\par
Q2 maximum $(-1,\,7)$; minimum $(2,\,-20)$\par
Q3 $k=-12$\par
Short solution:\par
Q1 $2x-8=0$; $y=16-32+3$; $y''=2>0$\par
Q2 $6x^{2}-6x-12=6(x-2)(x+1)$; $f''(x)=12x-6$: $f''(-1)=-18$, $f''(2)=18$\par
Q3 $3x^{2}+k=0$ at $x=2$: $12+k=0$\par
Watch for:\par
Q2 forgetting to divide out the $6$ is fine, but check both roots are found\par
Q3 is a Merit step: a turning point means $f'=0$, not $f=0$}
\end{frame}

% ------------------------------------------------------------------
\begin{frame}{Increasing and decreasing}\relax
\fbx{\text{increasing: } f'(x)>0 \qquad \text{decreasing: } f'(x)<0}
\vskip2mm
Find where $f'(x)=0$ first; test the intervals between.\\[1mm]
An up-parabola $f'(x)$ is \kw{negative between its roots}.
\vskip5mm
\qn{2} For $f(x)=x^{3}-6x^{2}+9x+1$, find where $f$ is decreasing.
\vskip4mm
\qn{3} For which $x$ is $y=5+4x-x^{2}$ increasing?
\sfoot{\hp{incdec}}
\note{Q2 $1<x<3$\par
Q3 $x<2$\par
Short solution:\par
Q2 $f'(x)=3(x-1)(x-3)$, up-parabola, negative between $1$ and $3$\par
Q3 $\dfrac{dy}{dx}=4-2x>0$\par
Watch for:\par
L3: an answer like “$x=1$ and $x=3$” is the turning points, not an interval\par
Q3 dividing by a negative flips the inequality if the student writes $-2x>-4$}
\end{frame}

% ------------------------------------------------------------------
\begin{frame}{Practice}\relax
\qn{1} For which values of $x$ is $f(x)=x^{2}-10x+1$ decreasing?
\vskip5mm
\qn{2} For which values of $x$ is $f(x)=x^{3}-12x+4$ decreasing?
\vskip5mm
\qn{3} Show that $f(x)=x^{3}+3x+1$ is increasing for all values of $x$.
\sfoot{\hp{incdec}}
\note{Q1 $x<5$\par
Q2 $-2<x<2$\par
Q3 $f'(x)=3x^{2}+3>0$ for all $x$\par
Short solution:\par
Q1 $f'(x)=2x-10<0$\par
Q2 $f'(x)=3x^{2}-12=3(x-2)(x+2)<0$ between the roots\par
Q3 $x^{2}\ge0$, so $3x^{2}+3\ge3>0$\par
Watch for:\par
Q2 writing $x<2$ only; the up-parabola is negative on a bounded interval\par
Q3 Merit: the reason ($x^{2}\ge0$) must be stated}
\end{frame}

% ------------------------------------------------------------------
\begin{frame}{Optimisation}\relax
{\sEmph\bfseries Write it in this order\par}\vskip1mm
\begin{enumerate}
\item Formula in \kw{one variable}.
\item Differentiate; set $=0$; solve.
\item Justify maximum or minimum.
\item Answer the question asked, with units.
\end{enumerate}
\vskip3mm
\qn{4} A ball's height is $h=2+19.6t-4.9t^{2}$ metres after $t$ seconds. Find the maximum height.
\sfoot{\hp{opt}}
\note{Q4 $21.6$ m\par
Short solution:\par
Q4 $\dfrac{dh}{dt}=19.6-9.8t=0$, $t=2$\par
Q4 $h=2+39.2-19.6$; $\dfrac{d^{2}h}{dt^{2}}=-9.8<0$\par
Watch for:\par
L5: answering $t=2$ s instead of the height\par
L6: the unit (m)}
\end{frame}

% ------------------------------------------------------------------
\begin{frame}{Worked example \textperiodcentered\ build the formula}\relax
\begin{columns}[T]
\begin{column}{0.56\textwidth}
\qn{5} A farmer has $60$ m of fence to make a rectangular pen against a wall, fencing three sides. Find the maximum area.
\vskip3mm
Step 1: area in terms of $x$ only.
\end{column}
\begin{column}{0.42\textwidth}\centering\vskip2mm
\setlength\figunit{11mm}
% Rectangular pen against a wall: two sides x, one side 60 - 2x.  Caller sets \figunit.
\begin{tikzpicture}[x=\figunit, y=\figunit, line cap=round]
  \fill[pattern=north east lines, pattern color=gray] (-0.4,2.2) rectangle (5.4,2.5);
  \draw[line width=0.8pt] (-0.4,2.2) -- (5.4,2.2);
  \node[above, fill=white, inner sep=1pt] at (2.5,2.5) {wall};
  \draw[line width=1.2pt] (0,2.2) -- (0,0) -- (5,0) -- (5,2.2);
  \node[left] at (0,1.1) {$x$};
  \node[right] at (5,1.1) {$x$};
  \node[below] at (2.5,0) {$60-2x$};
\end{tikzpicture}

\end{column}
\end{columns}
\sfoot{\hp{opt}}
\note{Q5 $450\ \mathrm{m}^{2}$ (pen $15$ m by $30$ m)\par
Short solution:\par
Q5 $A=x(60-2x)=60x-2x^{2}$\par
Q5 $\dfrac{dA}{dx}=60-4x=0$, $x=15$; $\dfrac{d^{2}A}{dx^{2}}=-4<0$\par
Watch for:\par
L4: writing $A=xy$ and differentiating with two letters\par
using $2x+2y=60$ (four sides) when the wall is one side}
\end{frame}

% ------------------------------------------------------------------
\begin{frame}{Practice}\relax
\begin{columns}[T]
\begin{column}{0.58\textwidth}
\qn{1} The profit, in dollars, from selling $x$ items is $P=120x-2x^{2}-500$. How many items give the maximum profit, and what is that profit?
\vskip5mm
\qn{2} An open box is made from a $12$ cm by $12$ cm sheet by cutting a square of side $x$ cm from each corner and folding up the sides. Find the value of $x$ that gives the maximum volume, and the maximum volume.
\end{column}
\begin{column}{0.38\textwidth}\centering\vskip18mm
\setlength\figunit{12mm}
% 12 cm square sheet with an x by x square cut from each corner; dashed fold lines.
% Drawn with x = 2.2 of 12 (scale only).  Caller sets \figunit (one unit = 1 cm of the sheet).
\begin{tikzpicture}[x=0.25\figunit, y=0.25\figunit, line cap=round]
  \def\c{2.2}
  \draw[line width=1pt] (\c,0) -- (12-\c,0) -- (12-\c,\c) -- (12,\c) -- (12,12-\c) -- (12-\c,12-\c)
     -- (12-\c,12) -- (\c,12) -- (\c,12-\c) -- (0,12-\c) -- (0,\c) -- (\c,\c) -- cycle;
  \draw[dashed, line width=0.6pt] (\c,\c) rectangle (12-\c,12-\c);
  \draw[gridcol, line width=0.5pt] (0,0) rectangle (\c,\c) (12-\c,0) rectangle (12,\c)
     (0,12-\c) rectangle (\c,12) (12-\c,12-\c) rectangle (12,12);
  \node[below, inner sep=2pt] at ({\c/2},0) {$x$};
  \node[left, inner sep=2pt] at (0,{\c/2}) {$x$};
  \draw[{Stealth[length=1.6mm]}-{Stealth[length=1.6mm]}, line width=0.5pt] (0,-2.6) -- (12,-2.6);
  \node[below, inner sep=2pt] at (6,-2.6) {$12$ cm};
\end{tikzpicture}

\end{column}
\end{columns}
\sfoot{\hp{opt}}
\note{Q1 $30$ items; profit 1300 dollars\par
Q2 $x=2$ cm; volume $128\ \mathrm{cm}^{3}$\par
Short solution:\par
Q1 $P'=120-4x=0$; $P=3600-1800-500$; $P''=-4<0$\par
Q2 $V=x(12-2x)^{2}=4x^{3}-48x^{2}+144x$\par
Q2 $V'=12x^{2}-96x+144=12(x-2)(x-6)$; $x=6$ gives no box\par
Q2 $V''=24x-96=-48<0$ at $x=2$\par
Watch for:\par
Q2 Excellence: the model must be built and the impossible root rejected with a reason}
\end{frame}

% ------------------------------------------------------------------
\begin{frame}{Ways to lose marks}\relax
\begin{itemize}
\item[\lmnum{noy}] A \tturn needs its $y$-value too.
\item[\lmnum{nojust}] ``Maximum'' needs a test and a sentence.
\item[\lmnum{interval}] Increasing or decreasing is an interval.
\item[\lmnum{twovars}] One variable before you differentiate.
\item[\lmnum{wrongq}] Answer the quantity asked, with units.
\end{itemize}
\sfoot{\hp{lose}}
\note{Teaching line, not a question slide.\par
Full list L1 to L6 in the handout; L6 (units and context) applies to every answer.}
\end{frame}

% ------------------------------------------------------------------
\begin{frame}{Summary}\relax
{\sSum
\begin{enumerate}
\item Turning point: $f'(x)=0$; then $y$ from $f(x)$.
\item Nature: $f''<0$ maximum, $f''>0$ minimum.
\item Increasing $f'>0$; decreasing $f'<0$: an interval.
\item Optimise: one variable, differentiate, $=0$, justify, answer.
\end{enumerate}}
\sfoot{\hp{idea}}
\note{Teaching line, not a question slide.\par
Ask the student to say each line back before looking at the screen.}
\end{frame}

% ------------------------------------------------------------------
\begin{frame}{Practice sheet \textperiodcentered\ 30 minutes}\relax
Work alone on the practice sheet.
\vskip3mm
\begin{itemize}
\item Questions 1 and 2: warm-up, about 6 minutes.
\item Questions 3 to 6: exam questions, about 20 minutes.
\item Question 7: Extension, if time allows.
\end{itemize}
\vskip3mm
Every maximum needs a reason; every answer needs units.
\sfoot{Practice sheet}
\note{Q1 maximum $(-3,\,64)$; minimum $(3,\,-44)$\par
Q2 $12.5$ m (at $t=1.5$ s)\par
Q3 $-4<x<2.5$\par
Q4 $26.8$ m\par
Q5 (i) $352$ and $-352$ viewers per month; (ii) $6336$ viewers\par
Q6 wall height $3$\par
Q7 (i) $169.8\ \mathrm{cm}^{2}$; (ii) $A''=-6\pi<0$\par
Short solution:\par
Q3 $f'(x)=2x^{2}+3x-20=(2x-5)(x+4)$, negative between the roots\par
Q4 $h'=22.5-9.8t=0$, $t=2.296$; $h''=-9.8<0$\par
Q5 (i) $t^{2}-48t+320=0$, $t=8$ or $40$; $V'=-22t+528$; (ii) $t=24$\par
Q6 $A=2x(9-x^{2})+x\cdot x^{2}=18x-x^{3}$; $x^{2}=6$; wall $9-6$\par
Q7 $L=40-\pi r$; $A=2rL-\pi r^{2}=80r-3\pi r^{2}$; $r=\dfrac{40}{3\pi}$\par
Watch for:\par
Q5 (i) two times give two rates, one positive, one negative\par
Q6 the triangle's height is $x^{2}$, not $9$\par
Full working: practice answer version (tutor's working, not an official mark scheme)}
\end{frame}

% ------------------------------------------------------------------
\begin{frame}{Homework}\relax
{\sEmph\bfseries Homework sheet: 14 questions, about 60 minutes\par}
\vskip4mm
\begin{itemize}
\item Questions 1 to 11: Achieved practice.
\item Questions 12 and 13: Merit.
\item Question 14: Extension (Excellence).
\end{itemize}
\vskip4mm
Optional extra: StudyTime Guide, ``Calculating Turning Points'' from page 16.
\sfoot{Homework}
\note{Teaching line, not a question slide.\par
Answers: homework answer version (tutor's working, not an official mark scheme).\par
Next lesson Do Now checks a turning point and a decreasing interval.}
\end{frame}
\end{document}
